\chapter{Antecedentes}
\label{cap:antecedentes}

\section{Los protocolos se apilan, no compiten}

Al enumerar «WiFi», «LoRaWAN» y «Contact ID» parece tratarse de alternativas entre las cuales elegir.
Operan en niveles distintos: la capa de \textbf{aplicación} define qué significa el mensaje (Contact
ID, JSON, SenML); la de \textbf{red y transporte}, cómo llega (IP, TCP/UDP, LoRaWAN, DC-09); la
\textbf{física}, cómo viajan los bits (WiFi 802.11, LoRa, cable). La independencia entre niveles es lo
que hace posible la propuesta: \textbf{el mismo mensaje Contact ID llega por radio LoRa, por WiFi o por
cable, y la plataforma lo interpreta de forma idéntica} (figura~\ref{fig:pila}).

\begin{figure}[H]
\centering
\resizebox{0.86\linewidth}{!}{%
\begin{tikzpicture}[font=\footnotesize]

\def\w{7.2}   % ancho de las cajas de capa
\def\h{1.15}  % alto de cada capa

% ---- capas ----
\node[caja, minimum width=\w cm, minimum height=\h cm, align=left,
      draw=azuloscuro, line width=1.1pt, fill=azuloscuro!5] (app) at (0,0)
  {\textbf{\color{azuloscuro}APLICACIÓN}\ \ \small ¿qué \emph{significa} el mensaje?\\[2pt]
   \scriptsize \textbf{SenML (RFC 8428)} · JSON · CSV del fabricante};

\node[caja, minimum width=\w cm, minimum height=\h cm, align=left,
      draw=verdeacento, line width=1.1pt, fill=verdeacento!5, below=1.5mm of app] (red)
  {\textbf{\color{verdeacento}RED Y TRANSPORTE}\ \ \small ¿cómo \emph{llega} de un punto a otro?\\[2pt]
   \scriptsize IP · TCP · UDP · \textbf{MQTT} · HTTP · \textbf{LoRaWAN}};

\node[caja, minimum width=\w cm, minimum height=\h cm, align=left,
      draw=ambar, line width=1.1pt, fill=ambar!5, below=1.5mm of red] (fis)
  {\textbf{\color{ambar}FÍSICA Y RADIO}\ \ \small ¿cómo \emph{viajan los bits}?\\[2pt]
   \scriptsize \textbf{WiFi (802.11)} · \textbf{LoRa} · Ethernet · 4G};

% ---- analogia de la carta ----
\node[caja, minimum width=5.4cm, minimum height=\h cm, align=left, right=8mm of app,
      draw=grislinea, fill=grisclaro] (a1)
  {\textbf{El idioma} en que está escrita\\[2pt]\scriptsize la carta};
\node[caja, minimum width=5.4cm, minimum height=\h cm, align=left, right=8mm of red,
      draw=grislinea, fill=grisclaro] (a2)
  {\textbf{Las direcciones} y el correo\\[2pt]\scriptsize que la encaminan};
\node[caja, minimum width=5.4cm, minimum height=\h cm, align=left, right=8mm of fis,
      draw=grislinea, fill=grisclaro] (a3)
  {\textbf{El camión, el avión}\\[2pt]\scriptsize que la transportan};

\node[nota, above=1.5mm of a1, font=\scriptsize\bfseries\color{gristexto}] {LA ANALOGÍA};

% ---- conclusion ----
\node[notaclave, below=6mm of fis.south, anchor=north, text width=13.4cm, xshift=3.1cm,
      fill=verdeacento!6, font=\footnotesize]
  {La misma carta en español puede viajar en camión, avión o bicicleta: \textbf{el idioma no cambia
   porque cambie el vehículo.} Del mismo modo, \textbf{la misma trama SenML llega por WiFi, por cable o
   por la rama LoRa, y la plataforma la interpreta igual en los tres casos.}};

\end{tikzpicture}}
\caption{Los protocolos se apilan, no compiten. Cada nivel responde una pregunta distinta, y la
elección en un nivel no condiciona la de los otros.}
\label{fig:pila}
\end{figure}


\section{Tecnologías de radio: LoRa y LoRaWAN}

\textbf{LoRa} es la técnica de modulación de radio (capa física, patente Semtech); \textbf{LoRaWAN} es
el protocolo de red abierto construido sobre ella: registro de dispositivos, cifrado AES-128 de extremo
a extremo y entrega al servidor \parencite{rp002}. La analogía exacta: LoRa es a LoRaWAN lo que 802.11
es a TCP/IP. Chile emplea la banda libre de 915--928\,MHz (Resolución 737 de Subtel
\parencite{lorachile}).

LoRa intercambia velocidad por alcance y autonomía: 0,3--27\,kbps y payloads de 11 a 53 bytes según la
tasa de datos, a cambio de 2--15\,km de alcance y años de batería. Su patrón de uso —lecturas breves y
avisos— coincide con el de los sensores de este trabajo. La cadena completa es nodo $\rightarrow$
gateway (traductor sin claves) $\rightarrow$ servidor de red ChirpStack (descifra, deduplica,
administra, reporta RSSI/SNR y publica en MQTT) $\rightarrow$ plataforma (figura~\ref{fig:lorawan}).

\begin{figure}[htbp]
\centering
\resizebox{\linewidth}{!}{%
\begin{tikzpicture}[font=\scriptsize]

\node[caja, text width=2.5cm, align=center, draw=ambar, line width=1pt] (nodos)
  {\textbf{NODOS}\\ sensores a batería\\ \emph{duermen y despiertan}};

\node[caja, text width=2.6cm, align=center, right=1.5cm of nodos, draw=ambar, line width=1pt] (gw)
  {\textbf{GATEWAY}\\ antena\\ \emph{traductor: no descifra}};

\node[caja, text width=3.0cm, align=center, right=1.5cm of gw, draw=verdeacento, line width=1.1pt,
      fill=verdeacento!5] (cs)
  {\textbf{ChirpStack}\\ servidor de red\\ \emph{descifra, deduplica, administra}};

\node[caja, text width=2.6cm, align=center, right=1.5cm of cs, draw=azuloscuro, line width=1.1pt,
      fill=azuloscuro!5] (plat)
  {\textbf{PLATAFORMA}\\ \emph{aporte de este trabajo}};

\draw[flecha] (nodos) -- node[nota, above, font=\scriptsize] {radio LoRa} node[nota, below, font=\scriptsize] {915 MHz} (gw);
\draw[flecha] (gw) -- node[nota, above, font=\scriptsize] {internet} node[nota, below, font=\scriptsize] {UDP/IP} (cs);
\draw[flechagruesa] (cs) -- node[nota, above, font=\scriptsize\bfseries\color{verdeacento}] {MQTT} (plat);

\node[notaaviso, below=5mm of gw.south, anchor=north, text width=13.2cm, xshift=3.2cm,
      fill=ambar!6, font=\scriptsize]
  {\textbf{El cifrado es de extremo a extremo, del nodo al servidor de red.} El gateway solo empaqueta
   tramas de radio y las reenvía: no posee las claves y no puede leer el contenido. Su única función es
   escuchar el aire.};

\end{tikzpicture}}
\caption{Cadena de una red LoRaWAN y ubicación del aporte de este trabajo.}
\label{fig:lorawan}
\end{figure}


\section{Mensajería MQTT}

MQTT aplica el patrón publicar/suscribir a través de un \emph{broker} (figura~\ref{fig:mqtt}). Dos
propiedades resultan determinantes: (1) \textbf{el dispositivo inicia la conexión saliente}, lo que
evita abrir puertos hacia los sensores y funciona tras NAT y cortafuegos —materializa el requerimiento
RNF-10—; (2) el mecanismo \emph{Last Will and Testament} notifica automáticamente la desconexión
abrupta de un cliente, lo que se traduce directamente al código 381 de Contact ID sin lógica adicional.

\begin{figure}[H]
\centering
\resizebox{0.8\linewidth}{!}{%
\begin{tikzpicture}[font=\scriptsize]

\node[caja, text width=1.9cm] (s1) at (0,1.35) {nodo Wi-Fi\\(ESP32)};
\node[caja, text width=1.9cm] (s2) at (0,0)    {pasarela LoRaWAN};
\node[caja, text width=1.9cm] (s3) at (0,-1.35){emulador de carga};

\node[caja, text width=2.5cm, align=center, draw=verdeacento, line width=1.2pt,
      fill=verdeacento!6, minimum height=2.9cm] (br) at (4.3,0)
  {\textbf{BROKER}\\[3pt] Mosquitto\\[3pt] \emph{recibe de quien publica,\\ entrega a quien se suscribe}};

\node[caja, text width=2.5cm] (c1) at (8.9,1.55) {el \textbf{backend}\\(seis filtros)};
\node[caja, text width=2.5cm] (c2) at (8.9,0)    {la base de datos\\(TimescaleDB)};
\node[caja, text width=2.5cm] (c3) at (8.9,-1.55){el \emph{dashboard}\\(navegador)};

\foreach \s in {s1,s2,s3} \draw[flecha] (\s.east) -- (br.west);
\draw[flecha, color=verdeacento] (br.east) -- (c1.west);
\draw[flecha] (c1.south) -- node[nota, left=0.5mm]{guarda} (c2.north);
\draw[flecha] (c1.east) -- ++(0.35,0) |- node[nota, pos=0.25, right=1mm]{HTTPS} (c3.east);

\node[nota, above=1mm of s1, font=\scriptsize\bfseries\color{gristexto}] {PUBLICAN};
\node[nota, above=1mm of c1, font=\scriptsize\bfseries\color{gristexto}] {SE SUSCRIBE};

\node[caja, below=8mm of br.south, anchor=north, text width=13.2cm, xshift=0.15cm,
      draw=grislinea, fill=grisclaro, font=\scriptsize, align=left]
  {Los mensajes se organizan en \textbf{tópicos}, con forma de ruta. En esta plataforma:
   \texttt{sitio/\textless sitio\textgreater/dispositivo/\textless id\_fisico\textgreater/\textless magnitud\textgreater}, por ejemplo \texttt{sitio/uct/dispositivo/dev-0001/temperatura}; las pasarelas, bajo \texttt{sitio/\textless sitio\textgreater/pasarela/\textless id\textgreater/…}\\[3pt]
   El backend es el \textbf{único suscriptor} (\texttt{sitio/\#}): recibe todo \textbf{sin necesidad de saber qué sensores
   existen}. Cada dispositivo solo puede publicar en sus propios tópicos, y el \emph{dashboard} recibe los datos del backend, no del broker.};

\end{tikzpicture}}
\caption{Patrón publicar/suscribir de MQTT en este sistema. El desacoplamiento entre emisores y receptor es lo que
permite sumar dispositivos sin tocar el backend.}
\label{fig:mqtt}
\end{figure}


\section{El protocolo Contact ID (SIA DC-05)}

Contact ID es el protocolo de reporte de eventos de la industria de alarmas, formalizado por la
Security Industry Association como estándar DC-05 \parencite{siadc05}. Tres estándares SIA se apilan:
\textbf{DC-05} define el formato del mensaje; DC-03 \parencite{siadc03} es un formato alternativo más
expresivo; \textbf{DC-09} \parencite{siadc09} transporta cualquiera de los dos sobre IP con cifrado
AES opcional y, en su revisión 2026, rotación de claves y auto-comisionamiento.

El mensaje es una secuencia de \textbf{16 dígitos} (figura~\ref{fig:mensaje}): cuenta (ACCT), tipo
(MT=18), calificador (\texttt{Q}: 1 abre una condición, 3 la cierra —una máquina de estados nativa—),
código de evento (XYZ), partición (GG), zona (CCC) y suma de comprobación módulo 15. El alfabeto tiene
15 símbolos porque el transporte original fue DTMF: tonos telefónicos, donde el cero vale diez.

\begin{figure}[htbp]
\centering
\resizebox{\linewidth}{!}{%
\begin{tikzpicture}

% ---- campos del mensaje ----
\node[campo, minimum width=17mm, fill=azuloscuro!8] (acct) {1234};
\node[campo, minimum width=11mm, right=0pt of acct, fill=grisclaro] (mt) {18};
\node[campo, minimum width=8mm,  right=0pt of mt,   fill=ambar!15]  (q)   {1};
\node[campo, minimum width=14mm, right=0pt of q,    fill=verdeacento!12] (xyz) {381};
\node[campo, minimum width=11mm, right=0pt of xyz,  fill=grisclaro] (gg)  {00};
\node[campo, minimum width=14mm, right=0pt of gg,   fill=verdeacento!12] (ccc) {007};
\node[campo, minimum width=8mm,  right=0pt of ccc,  fill=grisclaro] (s)   {B};

% ---- nombres de campo encima ----
\foreach \n/\l in {acct/ACCT, mt/MT, q/Q, xyz/XYZ, gg/GG, ccc/CCC, s/S}
  \node[font=\scriptsize\bfseries\color{azuloscuro}, above=1pt of \n] {\l};

% ---- descripciones debajo, escalonadas ----
\node[nota, below=17mm of s.south, xshift=6mm, anchor=west]   (ds)   {checksum};
\node[nota, below=24mm of ccc.south, xshift=6mm, anchor=west] (dccc) {zona 007: \textbf{el dispositivo}};
\node[nota, below=31mm of gg.south, xshift=6mm, anchor=west]  (dgg)  {grupo o partición};
\node[nota, below=38mm of xyz.south, xshift=6mm, anchor=west] (dxyz) {código de evento 381: \textbf{pérdida de supervisión RF}};
\node[nota, below=45mm of q.south, xshift=6mm, anchor=west]    (dq)   {calificador 1: \textbf{nuevo evento} (3 = restauración)};
\node[nota, below=52mm of mt.south, xshift=6mm, anchor=west]  (dmt)  {tipo de mensaje, siempre 18};
\node[nota, below=59mm of acct.south, xshift=6mm, anchor=west](dacct){cuenta 1234: \textbf{el sitio}};

\foreach \c/\d in {s/ds, ccc/dccc, gg/dgg, xyz/dxyz, q/dq, mt/dmt, acct/dacct} {
  \draw[flecha, dashed, line width=0.4pt] (\c.south) |- (\d.west);
}

% ---- lectura ----
\node[cajaacento, below=6mm of dacct.south west, anchor=north west, text width=12.5cm,
      fill=verdeacento!6, font=\footnotesize\itshape]
  {«En la cuenta 1234 se produjo un nuevo evento de pérdida de supervisión RF
    en el dispositivo 007.»};

\end{tikzpicture}}
\caption{Estructura del mensaje Contact ID: 16 dígitos distribuidos en siete campos.}
\label{fig:mensaje}
\end{figure}


La transmisión original incluía el \emph{kissoff}: la receptora confirma solo si la suma de
comprobación es válida, y el panel reintenta en su ausencia —confirmación a nivel de aplicación que
este trabajo conserva—. La «comunicación celular» de las alarmas modernas es un cambio de transporte,
no de protocolo (figura~\ref{fig:alarmas}); el mensaje de 16 dígitos nació para caber en tonos sobre
una línea analógica, \textbf{y esa restricción de hace cuatro décadas es la que hoy lo hace apropiado
para LoRa}, un canal igualmente estrecho.

\begin{figure}[htbp]
\centering
\begin{tikzpicture}[font=\scriptsize]

% ---- sensores y panel ----
\node[caja, text width=2.2cm, align=center, draw=ambar, line width=1pt] (sens) at (0,0)
  {\textbf{Sensores}\\ inalámbricos\\ a batería};

\node[caja, text width=2.2cm, align=center, draw=ambar, line width=1pt] (panel) at (3.5,0)
  {\textbf{Panel} o\\ comunicador};

\draw[flecha] (sens) -- node[nota, above, font=\scriptsize, align=center] {radio\\propietaria} (panel);

% ---- central receptora ----
\node[caja, text width=2.5cm, align=center, draw=azuloscuro, line width=1.1pt,
      fill=azuloscuro!5] (arc) at (12.6,0)
  {\textbf{Central}\\ \textbf{receptora}\\ (ARC / CMS)};

% ---- punto de bifurcacion: el panel elige transporte ----
\coordinate (bif) at (5.5,0);
\draw[line width=0.7pt, color=gristexto] (panel.east) -- (bif);
\draw[line width=0.7pt, color=gristexto] (5.5,1.55) -- (5.5,-1.55);

\coordinate (conv) at (11.1,0);
\draw[line width=0.7pt, color=gristexto] (11.1,1.55) -- (11.1,-1.55);
\draw[flecha] (conv) -- (arc.west);

% ---- las tres generaciones, con etiqueta centrada en el tramo ----
\draw[flecha] (5.5,1.55) -- (11.1,1.55);
\node[nota, above=0.6mm of {(8.3,1.55)}, font=\scriptsize]
  {línea telefónica fija: \textbf{DTMF}\ \ \textcolor{gristexto}{\itshape (histórico)}};

\draw[flecha] (5.5,0.45) -- (11.1,0.45);
\node[nota, above=0.6mm of {(8.3,0.45)}, font=\scriptsize]
  {red celular: \textbf{DTMF} sobre 2G/3G\ \ \textcolor{gristexto}{\itshape (en retirada)}};

\draw[flechagruesa] (5.5,-1.15) -- (11.1,-1.15);
\node[nota, above=0.6mm of {(8.3,-1.15)}, font=\scriptsize\color{verdeacento}]
  {\textbf{internet: SIA DC-09}\ \ \itshape (actual, el adoptado)};

% ---- nota final ----
\node[notaclave, text width=13.2cm, fill=verdeacento!6, font=\scriptsize, align=left]
      at (6.3,-3.1)
  {\textbf{En los tres casos el mensaje de aplicación es el mismo: Contact ID.} La «comunicación
   celular» es un cambio de transporte, no de protocolo: los comunicadores convierten la marcación de
   línea fija en marcación sobre red móvil conservando el mensaje intacto.\\[3pt]
   Por eso el mensaje tiene \textbf{16 dígitos y ni uno más}: debía caber en tonos telefónicos sobre una
   línea analógica. \textbf{Y esa restricción de hace cuatro décadas es la que hoy lo hace apropiado
   para LoRa}, un canal igualmente estrecho.};

\end{tikzpicture}
\caption{Cadena de comunicación de una alarma monitoreada. El formato del mensaje se conserva a través
de tres generaciones de transporte.}
\label{fig:alarmas}
\end{figure}


Su vocabulario incluye los eventos de salud de dispositivos inalámbricos que este trabajo necesita:
pérdida de supervisión RF (381), batería baja (384), interferencia (344), pruebas periódicas (602/603),
además de condiciones medibles (temperatura 158/159, agua 154, accesos 134/400)
\parencite{contactidcodes}.

\section{Persistencia de series de tiempo}

Los datos de sensores son voluminosos y siempre temporales: cincuenta sensores cada cinco minutos
generan cinco millones de registros anuales. Se emplea TimescaleDB, extensión de PostgreSQL que
conserva el modelo relacional y el SQL estándar, añadiendo particionado automático por tiempo,
agregación por intervalos y políticas de retención. La evaluación frente a alternativas está en la
sección~\ref{sec:infra}.
